\documentclass[hidelinks]{article}
\usepackage[a4paper, total={7in, 10in}]{geometry}
\usepackage[dvipsnames]{xcolor}
\usepackage{amsmath}
\usepackage{tikz}
\usepackage{tkz-euclide}
\usepackage[unicode]{hyperref}
\usepackage[all]{hypcap}
\usepackage{fancyhdr}
\usepackage{tocloft} % Add this package
\usepackage{xcolor}
\definecolor{darkgreen}{RGB}{0,100,0}
\usepackage{listings}
\lstset{
  language=R,
  basicstyle=\ttfamily\small,
  keywordstyle=\color{red},
  commentstyle=\color{darkgreen},
  stringstyle=\color{darkgreen},
  numbers=left,
  numberstyle=\tiny,
  stepnumber=1,
  numbersep=5pt,
  frame=single,
  breaklines=true,
  showstringspaces=false
}
\usepackage{amsmath}
\usepackage{amssymb}
\usepackage{pgfplots}
\pgfplotsset{compat=1.18}
\usepackage{amsmath}
\usepackage{float}
\usepackage{amsthm}
\usepackage[shortlabels]{enumitem}



% Remove dots from table of contents
\renewcommand{\cftdot}{}
\renewcommand{\cftsecdotsep}{\cftnodots}
\renewcommand{\cftsubsecdotsep}{\cftnodots}
\newcommand{\RNum}[1]{\uppercase\expandafter{\romannumeral #1\relax}}


\renewcommand{\contentsname}{Table of Contents}
\usetikzlibrary{angles,calc, decorations.pathreplacing}
\definecolor{carminered}{rgb}{1.0, 0.0, 0.22}
\definecolor{capri}{rgb}{0.0, 0.75, 1.0}
\definecolor{brightlavender}{rgb}{0.75, 0.58, 0.89}

\title{\textbf{STAT 545 HW 1}}
\author{Andres Efren Rocha Jayasinha}
\setcounter{secnumdepth}{0}
\date{Updated: \today} % suppresses the printed date

\begin{document}
\hypersetup{bookmarksnumbered=true,}
\pagecolor{white}
\color{black}
\maketitle
\begin{Large}
\tableofcontents
\end{Large}%
\pagebreak

\section{Assignment 1: Introduction to Markov Chains}

\subsection{Exercise 2.1}
A Markov chain has transition matrix
\[
P=
\begin{pmatrix}
0.1 & 0.3 & 0.6\\
0   & 0.4 & 0.6\\
0.3 & 0.2 & 0.5
\end{pmatrix},
\]
with initial distribution $\alpha=(0.2,\,0.3,\,0.5)$. Find the following:
\begin{enumerate}
  \item[(a)] $P(X_7=3 \mid X_6=2)$
  \item[(b)] $P(X_9=2 \mid X_1=2,\; X_5=1,\; X_7=3)$
  \item[(c)] $P(X_0=3 \mid X_1=1)$
  \item[(d)] $E(X_2)$
\end{enumerate}

\noindent {\color{darkgreen} \textbf{Stochastic Process.}
A \emph{stochastic process} is a collection of random variables
$\{X_t : t \in I\}$.
The set $I$ is called the \emph{index set} of the process.
The random variables are defined on a common \emph{state space} $S$.
}

\medskip

\noindent {\color{darkgreen} \textbf{Markov Chain.}
Let $S$ be a discrete set. A \emph{Markov chain} is a sequence of random variables
$X_0, X_1, \ldots$ taking values in $S$ with the property that
\[
P(X_{n+1} = j \mid X_0 = x_0, \ldots, X_{n-1} = x_{n-1}, X_n = i)
= P(X_{n+1} = j \mid X_n = i),
\]
for all $x_0, \ldots, x_{n-1}, i, j \in S$ and all $n \ge 0$.
The set $S$ is called the \emph{state space} of the Markov chain.


\smallskip

\noindent A Markov chain is said to be \emph{time-homogeneous} if the transition
probabilities do not depend on $n$, that is,
\[
P(X_{n+1} = j \mid X_n = i) = P(X_1 = j \mid X_0 = i),
\quad \text{for all } i,j \in S \text{ and } n \ge 0.
\]
}

\medskip

\noindent{\color{darkgreen} \textbf{Bayes.} Let $A_1, A_2, \ldots, A_k$ be events such that $P(A_i) > 0$, $i = 1, 2, \ldots, k$. 
Assume further that $A_1, A_2, \ldots, A_k$ form a partition of the sample space $\mathcal{C}$. 
Let $B$ be any event. Then
\[
P(A_j \mid B) = \frac{P(A_j) P(B \mid A_j)}{\sum_{i=1}^k P(A_i) P(B \mid A_i)}, \tag{1.4.3}
\]
}


\noindent{\color{darkgreen} \textbf{Law of Total Probability.}
Let $A_1,A_2,\dots,A_k$ be events such that $P(A_i)>0$ for $i=1,2,\dots,k$, and assume $A_1,\dots,A_k$ form a partition of the sample space $\mathcal{C}$. 
Then for any event $B$,
\[
P(B)=\sum_{i=1}^k P(A_i)P(B\mid A_i). 
\]
}

\noindent \textbf{Solution}

\paragraph{(a)}
For a time-homogeneous Markov chain,
\[
P(X_{n+1}=j \mid X_n=i)=p_{ij}.
\]
Therefore,
\[
P(X_7=3 \mid X_6=2)=p_{2,3}.
\]
From the transition matrix, the entry in row $2$ and column $3$ is $0.6$. Hence,
\[
\boxed{P(X_7=3 \mid X_6=2)=0.6.}
\]

\paragraph{(b)}
By the Markov property, conditioning on $X_7$ makes the earlier states $X_1$ and $X_5$ irrelevant. Thus,
\[
P(X_9=2 \mid X_1=2,\; X_5=1,\; X_7=3)
= P(X_9=2 \mid X_7=3).
\]

\noindent Using the law of total probability and conditioning on the possible values of $X_8$,
\[
P(X_9=2 \mid X_7=3)
= \sum_{j=1}^3 P(X_9=2 \mid X_8=j,\; X_7=3)\,P(X_8=j \mid X_7=3).
\]

\noindent By the Markov property,

\[
P(X_9=2 \mid X_8=j,\; X_7=3)=P(X_9=2 \mid X_8=j)=p_{j2},
\]

\noindent and

\[
P(X_8=j \mid X_7=3)=p_{3j}.
\]

\noindent Therefore,

\[
P(X_9=2 \mid X_7=3)
= \sum_{j=1}^3 p_{3j}p_{j2}.
\]

\noindent Substituting the values from the transition matrix (or multiply second column by third row),

\[
=0.3(0.3)+0.2(0.4)+0.5(0.2)
=0.09+0.08+0.10
=0.27.
\]

\noindent Hence,

\[
\boxed{P(X_9=2 \mid X_1=2,\; X_5=1,\; X_7=3)=0.27.}
\]

\paragraph{(c)}
We want $P(X_0=3 \mid X_1=1)$. By Bayes' rule,
\[
P(X_0=3 \mid X_1=1)=\frac{P(X_1=1 \mid X_0=3)\,P(X_0=3)}{P(X_1=1)}.
\]
From the initial distribution $\alpha=(0.2,0.3,0.5)$, we have $P(X_0=3)=0.5$.
Also, from the transition matrix,
\[
P(X_1=1 \mid X_0=3)=p_{3,1}=0.3.
\]
Next,
\[
P(X_1=1)=\sum_{i=1}^3 P(X_1=1 \mid X_0=i)\,P(X_0=i)
=\sum_{i=1}^3 p_{i1}\,\alpha_i.
\]
Using column 1 of $P$, $(p_{11},p_{21},p_{31})=(0.1,0,0.3)$, we get
\[
P(X_1=1)=0.1(0.2)+0(0.3)+0.3(0.5)=0.02+0+0.15=0.17.
\]
Therefore,
\[
P(X_0=3 \mid X_1=1)=\frac{0.3\cdot 0.5}{0.17}=\frac{0.15}{0.17}
=\frac{15}{17}\approx 0.88235.
\]
\[
\boxed{P(X_0=3 \mid X_1=1)=\frac{15}{17}\approx 0.88235.}
\]

\paragraph{(d)}

To compute $E(X_2)$, first find the distribution of $X_2$.
Since the initial distribution is $\alpha=(0.2,0.3,0.5)$, we have
\[
\mathbb{P}(X_2=\cdot)=\alpha P^2.
\]

\noindent Compute

\[
P^2=
\begin{pmatrix}
0.19 & 0.27 & 0.54\\
0.18 & 0.28 & 0.54\\
0.18 & 0.27 & 0.55
\end{pmatrix}.
\]

\noindent Therefore,

\[
\alpha P^2
=
(0.2,0.3,0.5)
\begin{pmatrix}
0.19 & 0.27 & 0.54\\
0.18 & 0.28 & 0.54\\
0.18 & 0.27 & 0.55
\end{pmatrix}
=
(0.182,\;0.273,\;0.545).
\]

\noindent So

\[
\mathbb{P}(X_2=1)=0.182,\qquad \mathbb{P}(X_2=2)=0.273,\qquad \mathbb{P}(X_2=3)=0.545.
\]

\noindent Finally,

\[
E(X_2)=\sum_{k=1}^3 k\,\mathbb{P}(X_2=k)
=1(0.182)+2(0.273)+3(0.545)
=0.182+0.546+1.635
=2.363.
\]
\[
\boxed{E(X_2)=2.363.}
\]

\subsection{Exercise 2.9} 

Give the transition matrix for the transition graph in Figure 2.11

\begin{tikzpicture}[
  >=Stealth,
  node/.style={circle, draw, thick, minimum size=9mm, inner sep=0pt, font=\itshape},
  lab/.style={font=\small},
  edge/.style={->, thick}
]

% Nodes (hand-placed to resemble the figure)
\node[node] (a) at (0,0) {$a$};
\node[node] (b) at (2,2.3) {$b$};
\node[node] (c) at (4,0) {$c$};
\node[node] (e) at (1,-2.0) {$e$};
\node[node] (d) at (3,-2.0) {$d$};

% Edges
\draw[edge] (b) -- node[lab, midway, above left] {$1$} (a);

\draw[edge] (a) -- node[lab, midway, above] {$3$} (c);

\draw[edge] (c) -- node[lab, midway, above right] {$2$} (b);

\draw[edge] (c) -- node[lab, midway, right] {$1$} (d);

\draw[edge] (b) -- node[lab, midway, left] {$4$} (e);

\draw[edge] (d) -- node[lab, midway, right] {$5$} (b);

\draw[edge] (e) -- node[lab, midway, below] {$2$} (d);

% Curved pair between a and e
\draw[edge, bend left=25] (a) to node[lab, midway, left] {$2$} (e);
\draw[edge, bend left=25] (e) to node[lab, midway, left] {$6$} (a);

% Self-loops
\draw[edge] (b) edge[loop above] node[lab] {$2$} (b);
\draw[edge] (c) edge[loop right] node[lab] {$6$} (c);

\end{tikzpicture}

\noindent \textbf{Solution}

% States ordered as (a,b,c,d,e).
% Read weights from the graph:
% a->c:3, a->e:2
% b->a:1, b->b:2, b->e:4
% c->b:2, c->c:6, c->d:1
% d->b:5
% e->a:6, e->d:2

\[
\textbf{Outgoing totals:}\quad
\begin{aligned}
&\deg^+(a)=3+2=5, \qquad
\deg^+(b)=1+2+4=7,\\
&\deg^+(c)=2+6+1=9, \qquad
\deg^+(d)=5, \qquad
\deg^+(e)=6+2=8.
\end{aligned}
\]


\[
\textbf{Transition probabilities:}\quad
P_{ij}=\frac{w(i\to j)}{\sum_{k} w(i\to k)}.
\]

\[
\implies \textbf{Transition matrix }P\text{ (order }(a,b,c,d,e)\text{):}\quad
P=
\begin{pmatrix}
0 & 0 & \frac{3}{5} & 0 & \frac{2}{5} \\[6pt]
\frac{1}{7} & \frac{2}{7} & 0 & 0 & \frac{4}{7} \\[6pt]
0 & \frac{2}{9} & \frac{2}{3} & \frac{1}{9} & 0 \\[6pt]
0 & 1 & 0 & 0 & 0 \\[6pt]
\frac{3}{4} & 0 & 0 & \frac{1}{4} & 0
\end{pmatrix}
\]

\subsection{Exercise 2.10} Consider a Markov chain with transition matrix
\[
P=
\begin{array}{c|cccc}
 & a & b & c & d\\ \hline
a & 0   & \frac{3}{5} & \frac{1}{5} & \frac{1}{5}\\
b & \frac{3}{4} & 0   & \frac{1}{4} & 0\\
c & \frac{1}{4} & \frac{1}{4} & \frac{1}{4} & \frac{1}{4}\\
d & \frac{1}{4} & 0   & \frac{1}{4} & \frac{1}{2}
\end{array}
\]

\begin{enumerate}
  \item[(a)] Exhibit the directed, weighted transition graph for the chain.
  \item[(b)] The transition graph for this chain can be given as a weighted graph without
  directed edges. Exhibit the graph.
\end{enumerate}

\noindent \textbf{Solution}

\paragraph{(a)}

\begin{center}
\begin{tikzpicture}[
  >=Stealth,
  node/.style={circle, draw, thick, minimum size=10mm, inner sep=0pt, font=\itshape},
  lab/.style={font=\small},
  edge/.style={->, thick}
]

% Nodes (spread out)
\node[node] (a) at (0,4) {$a$};
\node[node] (b) at (-4,0) {$b$};
\node[node] (c) at (0,0) {$c$};
\node[node] (d) at (4,0) {$d$};

% Edges from a
\draw[edge] (a) to[bend left=15] node[lab, midway, left] {$\frac{3}{5}$} (b);
\draw[edge] (a) -- node[lab, midway, right] {$\frac{1}{5}$} (c);
\draw[edge] (a) to[bend left=15] node[lab, midway, right] {$\frac{1}{5}$} (d);

% Edges from b
\draw[edge] (b) to[bend left=20] node[lab, midway, above] {$\frac{3}{4}$} (a);
\draw[edge] (b) -- node[lab, midway, above] {$\frac{1}{4}$} (c);

% Edges from c
\draw[edge] (c) -- node[lab, midway, above] {$\frac{1}{4}$} (b);
\draw[edge] (c) -- node[lab, midway, above] {$\frac{1}{4}$} (d);
\draw[edge] (c) to[bend left=25] node[lab, midway, left] {$\frac{1}{4}$} (a);
\draw[edge] (c) edge[loop below] node[lab] {$\frac{1}{4}$} (c);

% Edges from d
\draw[edge] (d) to[bend left=25] node[lab, midway, above] {$\frac{1}{4}$} (a);
\draw[edge] (d) -- node[lab, midway, below] {} (c);
\draw[edge] (d) edge[loop right] node[lab] {$\frac{1}{2}$} (d);

\end{tikzpicture}
\end{center}

\paragraph{(b)}

\begin{center}
\begin{tikzpicture}[
  node/.style={circle, draw, thick, minimum size=9mm, inner sep=0pt},
  lab/.style={font=\small}
]

% Nodes
\node[node] (a) at (0,3) {$a$};
\node[node] (b) at (-2,1) {$b$};
\node[node] (c) at (0,-1) {$c$};
\node[node] (d) at (2,1) {$d$};

% Undirected edges
\draw (a) -- node[lab, midway, left] {$3$} (b);
\draw (b) -- node[lab, midway, left] {$1$} (c);
\draw (c) -- node[lab, midway, right] {$1$} (d);
\draw (d) -- node[lab, midway, right] {$1$} (a);
\draw (a) -- node[lab, midway, right] {$1$} (c);

% Loops
\draw (c) edge[loop below] node[lab] {$1$} (c);
\draw (d) edge[loop right] node[lab] {$2$} (d);

\end{tikzpicture}
\end{center}

\subsection{Exercise 2.16} Assume that \(P\) is a stochastic matrix with equal rows.
Show that

\[
P^{n} = P,
\qquad \text{for all } n \ge 1.
\]

\noindent{\color{darkgreen} \textbf{(Stochastic Matrix).}
A stochastic matrix is a square matrix \(P\) such that:
\begin{enumerate}
  \item \(P_{ij} \ge 0\) for all \(i,j\).
  \item For each row \(i\),
  \[
  \sum_{j} P_{ij} = 1.
  \]
\end{enumerate}
}

\noindent \textbf{Solution}

\smallskip

\noindent Assume that \(P\) is a stochastic matrix whose rows are all equal. Let
\(r^{T}\) denote this common row, so that \(P_{ij}=r_j\) for all \(i,j\).
Since \(P\) is stochastic, the entries of \(r\) satisfy
\[
\sum_{k} r_k = 1.
\]

\noindent We will prove by induction for \(n\ge 1\) that \(P^{n}=P\).

\medskip

\noindent \textbf{Base case (\(n=1\)):}
Clearly,
\[
P^{1}=P.
\]

\medskip

\noindent \textbf{Inductive hypothesis:}
Assume that for some \(k\ge 1\),
\[
P^{k}=P.
\]

\medskip

\noindent \textbf{Inductive step:}
Using the inductive hypothesis,
\[
P^{k+1}=P^{k}P=PP=P^{2}.
\]
If we  now show that \(P^{2}=P\), then the proof is complete. For any \(i,j\),
\[
(P^{2})_{ij}=\sum_{m} P_{im}P_{mj}
=\sum_{m} r_m r_j
=r_j \sum_{m} r_m
=r_j
=P_{ij},
\]
where we used the fact that \(\sum_{m} r_m=1\).
Thus \(P^{2}=P\), and hence
\[
P^{k+1}=P.
\]

\medskip

\noindent By mathematical induction,
\[
P^{n}=P \qquad \text{for all } n\ge 1.
\]

\subsection{Exercise 2.27}

\textbf{R:} See \texttt{gamblersruin.R}. Simulate gambler's ruin for a gambler
with initial stake \$2, playing a fair game.
\begin{enumerate}
  \item[(a)] Estimate the probability that the gambler is ruined before he wins \$5.
  \item[(b)] Construct the transition matrix for the associated Markov chain.
  Estimate the desired probability in (a) by taking high matrix powers.
  \item[(c)] Compare your results with the exact probability.
\end{enumerate}

\noindent \textbf{Solution}

\paragraph{(a)}\mbox{}\\
\noindent\textbf{Input:}

\begin{lstlisting}[caption={Gambler's Ruin Simulation in R}]
# gamblersruin.R
# Example 1.11

# gamble(k, n, p)
#   k: Gambler's initial state
#   n: Gambler plays until either n or Ruin
#   p: Probability of winning $1 at each play
#   Function returns 1 if gambler is eventually ruined
#                    returns 0 if gambler eventually wins n

gamble <- function(k, n, p) {
  stake <- k
  while (stake > 0 & stake < n) {
    bet <- sample(c(-1, 1), 1, prob = c(1 - p, p))
    stake <- stake + bet
  }
  if (stake == 0) return(1) else return(0)
}

k <- 2
n <- 40
p <- 1/2
trials <- 1000

simlist <- replicate(trials, gamble(k, n, p))
mean(simlist)  # Estimate of probability that gambler is ruined

# For p = 0.5, exact probability is (n - k) / n
\end{lstlisting}


\noindent\textbf{Output:} \textit{Estimated {Probability of Ruin}} = \(\boxed{0.952}\)

\paragraph{(b)}\mbox{}\\
\noindent\textbf{Input:}

\begin{lstlisting}[caption={Gambler's Ruin Simulation in R}]
k <- 2
n <- 40
p <- 1/2
q <- 1 - p

P <- matrix(0, n+1, n+1)

# absorbing states
P[1,1] <- 1           # state 0
P[n+1, n+1] <- 1      # state n

# interior transitions
for (i in 1:(n-1)) {
  P[i+1, i]   <- q    # from i to i-1
  P[i+1, i+2] <- p    # from i to i+1
}

# high power
t <- 2000
Pt <- diag(n+1)
for (j in 1:t) Pt <- Pt %*% P

ruin_prob_est <- Pt[k+1, 1]     # (P^t)_{k,0}
win_prob_est  <- Pt[k+1, n+1]   # (P^t)_{k,n}

ruin_prob_est
win_prob_est
ruin_prob_est + win_prob_est    # should be ~1
\end{lstlisting}


\noindent\textbf{Output:} \textit{Estimated {Probability of Ruin}} = \(\boxed{0.949}\)

\smallskip

\noindent \textit{\textbf{Note that the transition matrix P is,}}

\[
P =
\begin{pmatrix}
1      & 0      & 0      & 0      & \cdots & 0      & 0 \\
q      & 0      & p      & 0      & \cdots & 0      & 0 \\
0      & q      & 0      & p      & \cdots & 0      & 0 \\
\vdots & \ddots & \ddots & \ddots & \ddots & \vdots & \vdots \\
0      & \cdots & 0      & q      & 0      & p      & 0 \\
0      & \cdots & 0      & 0      & q      & 0      & p \\
0      & 0      & 0      & 0      & \cdots & 0      & 1
\end{pmatrix}
\]

\noindent \textit{\textbf{with p = $\frac{1}{2}$} and q = $\frac{1}{2}$}

\paragraph{(c)} From \textbf{gamblersruin.R} we have,

\[
\textit{Exact Probability of Ruin} = \frac{n-k}{n}
\]

\[
\implies \textit{Exact Probability of Ruin} = \frac{40-2}{40} = \boxed{0.95}
\]

\noindent Given that the estimated probabilities of ruin were 0.952, 0.949, they match up nicely with the exact probability of ruin, that being 0.95.

\end{document}